\expandafter\ifx\csname CAPETINCLUDEMODE\endcsname\relax
\documentclass[11pt,a4paper]{article}
% preamble.tex — packages et commandes communes a toutes les fiches CAPET SII IE
\input{glyphtounicode}
\pdfgentounicode=1
\usepackage{cmap}
\usepackage[T1]{fontenc}
\usepackage[utf8]{inputenc}
\usepackage[french]{babel}
\usepackage[a4paper, top=1.6cm, bottom=1.6cm, left=1.4cm, right=1.4cm]{geometry}
\setlength{\headheight}{14.5pt}
\usepackage{amsmath, amssymb, amsfonts}
\usepackage{siunitx}
\sisetup{locale=FR, output-decimal-marker={,}, inter-unit-product=\ensuremath{{}\cdot{}}}
\usepackage{array, booktabs, multirow, tabularx}
\usepackage{xcolor}
\usepackage[most, breakable]{tcolorbox}
\usepackage{multicol}
\usepackage{enumitem}
\usepackage{mathtools}
\usepackage{lmodern}
\usepackage{microtype}
\usepackage{fancyhdr}
\usepackage{tikz}
\usetikzlibrary{positioning,arrows.meta}
\usepackage{lastpage}

% ── Couleurs par domaine ──────────────────────────────────────────────
\definecolor{colElec}{RGB}{26, 82, 148}       % Électrotechnique — bleu foncé
\definecolor{colElecLight}{RGB}{210, 225, 245} % fond bleu clair
\definecolor{colAuto}{RGB}{34, 120, 66}        % Automatique — vert
\definecolor{colMeca}{RGB}{185, 90, 20}        % Mécanique — orange
\definecolor{colTherm}{RGB}{160, 40, 40}       % Thermique — rouge
\definecolor{colQT}{RGB}{232, 244, 233}        % fond question type (vert très léger)
\definecolor{colJury}{RGB}{255, 243, 205}      % fond jury (jaune ambre)
\definecolor{colJuryFrame}{RGB}{200, 130, 0}   % bordure jury
\definecolor{colGrayLight}{RGB}{245, 245, 247} % fond gris neutre
\definecolor{colSeparator}{RGB}{180, 180, 180} % gris séparateur

% ── Style tcolorbox de section principale ────────────────────────────
\tcbset{
  ficheSection/.style={
    breakable,
    enhanced,
    colback=#1!7!white,
    colframe=#1,
    colbacktitle=#1,
    coltitle=white,
    fonttitle=\bfseries\normalsize,
    left=4pt, right=4pt, top=3pt, bottom=4pt,
    before skip=6pt, after skip=6pt,
  }
}

% ── Style encadré question type ──────────────────────────────────────
\tcbset{
  questiontype/.style={
    breakable,
    enhanced,
    colback=colQT,
    colframe=colElec!55!black,
    colbacktitle=colElec!55!black,
    coltitle=white,
    fonttitle=\bfseries\small,
    left=4pt, right=4pt, top=3pt, bottom=3pt,
    before skip=5pt, after skip=5pt,
  }
}

% ── Style encadré jury ───────────────────────────────────────────────
\tcbset{
  jury/.style={
    breakable,
    enhanced,
    colback=colJury,
    colframe=colJuryFrame,
    fonttitle=\bfseries\normalsize,
    left=4pt, right=4pt, top=3pt, bottom=3pt,
  }
}

% ── Mise en page ─────────────────────────────────────────────────────
\pagestyle{fancy}
\fancyhf{}
\renewcommand{\headrulewidth}{0.4pt}
\renewcommand{\footrulewidth}{0.3pt}

\setlist[itemize]{leftmargin=1.2em, itemsep=1pt, topsep=2pt, parsep=0pt}
\setlist[enumerate]{leftmargin=1.2em, itemsep=3pt, topsep=2pt, parsep=0pt}

\setlength{\parskip}{2pt}
\setlength{\columnsep}{10pt}

\def\CAPETBEGINFILE{\begin{document}}
\def\CAPETENDFILE{\end{document}}
\else
\def\CAPETBEGINFILE{}
\def\CAPETENDFILE{}
\fi

\fancyhead[L]{\textbf{\textcolor{colElec}{T19}} - AOP (amplificateur operationnel)}
\fancyhead[R]{\small\textit{CAPET SII IE - Session 2026}}
\fancyfoot[C]{\small\thepage\,/\,\pageref{LastPage}}

\CAPETBEGINFILE

% -----------------------------------------------------------------------
%  BLOC EN-TeTE
% -----------------------------------------------------------------------

\begin{tcolorbox}[
  enhanced, breakable,
  colback=colElec,
  colframe=colElec,
  left=6pt, right=6pt, top=5pt, bottom=5pt,
  before skip=0pt, after skip=4pt
]
  {\color{white}\bfseries\large T19 \quad AOP (amplificateur operationnel)}\\[2pt]
  {\color{colElec!30!white}\small
    \textbf{Domaine :} electronique - Conditionnement du signal \quad
    \textbf{Presence :} \textbf{8/13} sujets (2013-2026) - theme \emph{tres frequent}
  }\\[1pt]
  {\color{colElec!20!white}\footnotesize
    Sous-themes : montages inverseur - non-inverseur - suiveur - soustracteur - integrateur - comparateur - hysteresis - conversion tension-courant 4-20\,mA - theoreme de Millman
  }
\end{tcolorbox}

% -----------------------------------------------------------------------
%  1. RAPPELS DE COURS
% -----------------------------------------------------------------------

\begin{tcolorbox}[
  ficheSection=colElec,
  title={1. Rappels de cours},
]

% -- 1.1 Hypotheses AOP ideal -----------------------------------------
\textbf{1.1 Hypotheses de l'AOP ideal - regime lineaire}

\smallskip
\begin{center}
\renewcommand{\arraystretch}{1.3}
\begin{tabular}{ll}
\toprule
\textbf{Hypothese} & \textbf{Consequence} \\
\midrule
Courants d'entree nuls & $I_+ = I_- = 0$ \\
Gain differentiel infini & $\varepsilon = V_+ - V_- = 0 \;\Rightarrow\; V_+ = V_-$ \\
\bottomrule
\end{tabular}
\end{center}
\textbf{Condition de validite :} contre-reaction negative (sortie rebouclee sur $V_-$). En comparateur (boucle ouverte ou contre-reaction positive), $V_+ \neq V_-$.

\medskip\hrule\vspace{4pt}

% -- 1.2 Montages fondamentaux ----------------------------------------
\textbf{1.2 Montages fondamentaux}

\smallskip
\begin{center}
\renewcommand{\arraystretch}{1.45}
\begin{tabular}{llll}
\toprule
\textbf{Montage} & \textbf{Gain $V_s/V_e$} & \textbf{$Z_e$} & \textbf{Usage} \\
\midrule
Inverseur & $-R_2/R_1$ & $R_1$ & Amplification + inversion \\
Non-inverseur & $1 + R_2/R_1$ & $\to \infty$ & Amplification sans inversion \\
Suiveur (buffer) & $1$ & $\to \infty$ & Adaptation d'impedance \\
Soustracteur & $\dfrac{R_2}{R_1}(V_2 - V_1)$ & - & Signal differentiel (capteur) \\
Integrateur & $-\dfrac{1}{RCp}$ & $R_1$ & Filtre passe-bas, correcteur I \\
\bottomrule
\end{tabular}
\end{center}

\medskip\hrule\vspace{4pt}

% -- 1.3 Comparateur - hysteresis -------------------------------------
\textbf{1.3 Comparateur - hysteresis - contre-reaction positive}

\smallskip
Contre-reaction sur $V_+$ (resistances $R_1$ sur sortie, $R_2$ sur $V_{ref}$) :
\[
  V_+ = \frac{R_2\,V_{ref} + R_1\,V_s}{R_1 + R_2}
\]
Deux seuils de basculement (avec $V_s = \pm V_{sat}$) :
\[
  V_H = \frac{R_2\,V_{ref} + R_1\,V_{sat}}{R_1 + R_2}
  \qquad;\qquad
  V_B = \frac{R_2\,V_{ref} - R_1\,V_{sat}}{R_1 + R_2}
\]

\medskip\hrule\vspace{4pt}

% -- 1.4 Conversion tension-courant -----------------------------------
\textbf{1.4 Conversion tensionecourant 4-20\,mA (transmetteur normalise)}

\smallskip
\[
  I_s = \frac{V_{entree}}{R_5}
\]
Pour $V_{entree} \in [1\,;\,5]\,\mathrm{V}$ et $R_5 = 250\,\Omega$ : $I_s \in [4\,;\,20]\,\mathrm{mA}$.\\
$I_s$ est independant de la charge $R_{ana}$ car la boucle de contre-reaction maintient $\varepsilon = 0$.

\textbf{Theoreme de Millman} (pour un neud - $n$ admittances) :
\[
  V_{neud} = \frac{\sum_k V_k/R_k}{\sum_k 1/R_k}
\]

\end{tcolorbox}

% -----------------------------------------------------------------------
%  2. FORMULAIRE SYNTHeTIQUE
% -----------------------------------------------------------------------

\begin{tcolorbox}[
  ficheSection=colElec,
  title={2. Formulaire synthetique},
]
\small
\begin{multicols}{2}
\renewcommand{\arraystretch}{1.45}
\begin{tabular}{@{}p{4.5cm}ll@{}}
\toprule
\textbf{Formule} & \textbf{Usage} & \textbf{U.} \\
\midrule
\multicolumn{3}{@{}l}{\footnotesize\textit{Montages lineaires}} \\
$V_s/V_e = -R_2/R_1$ & Inverseur & - \\
$V_s/V_e = 1+R_2/R_1$ & Non-inv. & - \\
$V_s = V_e$ & Suiveur & - \\
$V_s = (R_2/R_1)(V_2-V_1)$ & Soustr. & - \\
$H(p) = -1/(RCp)$ & Integr. & - \\
\midrule
\multicolumn{3}{@{}l}{\footnotesize\textit{Hypotheses}} \\
$V_+ = V_-$ & Lin. & - \\
$I_+ = I_- = 0$ & Lin. & - \\
\bottomrule
\end{tabular}

\columnbreak

\renewcommand{\arraystretch}{1.45}
\begin{tabular}{@{}p{4.5cm}ll@{}}
\toprule
\textbf{Formule} & \textbf{Usage} & \textbf{U.} \\
\midrule
\multicolumn{3}{@{}l}{\footnotesize\textit{Hysteresis}} \\
$V_+ = (R_2 V_{ref}+R_1 V_s)/(R_1+R_2)$ & Seuil & V \\
$V_H = (R_2 V_{ref}+R_1 V_{sat})/(R_1+R_2)$ & Haut & V \\
$V_B = (R_2 V_{ref}-R_1 V_{sat})/(R_1+R_2)$ & Bas & V \\
\midrule
\multicolumn{3}{@{}l}{\footnotesize\textit{T/I et Millman}} \\
$I_s = V_e/R_5$ & Transm. & A \\
$V_n = \sum(V_k/R_k)/\sum(1/R_k)$ & Millman & V \\
\midrule
\multicolumn{3}{@{}l}{\footnotesize\textit{Donnees 2024}} \\
$G_{total} = 340$ & Transmetteur & - \\
Plage : 4-20\,mA, $R_5 = 250\,\Omega$ & Normalise & - \\
\bottomrule
\end{tabular}
\end{multicols}
\normalsize
\end{tcolorbox}

% -----------------------------------------------------------------------
%  3. QUESTIONS TYPES
% -----------------------------------------------------------------------

\begin{tcolorbox}[
  ficheSection=colElec,
  title={3. Questions types - extraites des sujets reels},
]

% -- QT1 -------------------------------------------------------------

% Mini-procedure de resolution type copie concours
\begin{tcolorbox}[questiontype,title={Procedure type copie concours (5 etapes)}]
\begin{enumerate}[leftmargin=1.5em]
  \item \textbf{Qualifier le modele} : hypotheses, domaine de validite, grandeurs d'entree/sortie.
  \item \textbf{Ecrire les lois} : equations de base avec symboles et unites SI.
  \item \textbf{Isoler l'inconnue} : expression litterale, puis application numerique tracee.
  \item \textbf{Controler l'ordre de grandeur} : coherence physique et unite finale.
  \item \textbf{Conclure concours} : phrase courte orientee decision technique.
\end{enumerate}
\end{tcolorbox}

% Visuel de methode : chemin de resolution CAPET
\begin{center}
\begin{tikzpicture}[font=\small, node distance=7mm and 6mm, >=Stealth]
  \node[draw, rounded corners, fill=colElec!8, inner sep=4pt] (n1) {Notions};
  \node[draw, rounded corners, fill=colElec!8, right=of n1, inner sep=4pt] (n2) {Methode};
  \node[draw, rounded corners, fill=colElec!8, right=of n2, inner sep=4pt] (n3) {Application annale};
  \node[draw, rounded corners, fill=colElec!8, right=of n3, inner sep=4pt] (n4) {Verification jury};
  \draw[->, thick, colElec] (n1) -- (n2);
  \draw[->, thick, colElec] (n2) -- (n3);
  \draw[->, thick, colElec] (n3) -- (n4);
\end{tikzpicture}
\end{center}
\begin{tcolorbox}[questiontype,
  title={QT1 - Montage differentiel - calcul du gain $G_1$ (transmetteur de temperature) \hfill\normalfont\footnotesize E2 disc. 2024 (E1-E3)}
]
\textbf{enonce.}
Un montage AOP differentiel (soustracteur) amplifie la tension $e_{AB}$ d'un capteur PT100. Donnees : $R_1 = 1\,\mathrm{k\Omega}$, $R_2 = 834{,}4\,\Omega$. Exprimer le gain $G_1$ et calculer sa valeur. Le 2\textsuperscript{e} etage a $G_2 = -100$, le 3\textsuperscript{e} a $G_3 = -1$. Calculer le gain global.

\tcbline

\textbf{Corrige.}

Par la loi des mailles (AOP ideal, $V_+ = V_-$) :
\[
  G_1 = \frac{2R_1 + R_2}{R_2} = \frac{2 \times 1000 + 834{,}4}{834{,}4} = \frac{2834{,}4}{834{,}4} \approx \mathbf{3{,}4}
\]
Gain global :
\[
  G = G_1 \times G_2 \times G_3 = 3{,}4 \times (-100) \times (-1) = \mathbf{340}
\]
$V_3 = 340 \times e_{AB}$ - les deux inversions de $G_2$ et $G_3$ se compensent, le signal reste en phase avec $e_{AB}$.
\end{tcolorbox}

\vspace{4pt}

% -- QT2 -------------------------------------------------------------

\begin{tcolorbox}[questiontype,
  title={QT2 - Comparateur - hysteresis - seuils de basculement \hfill\normalfont\footnotesize E2 disc. 2023 (F6-F8)}
]
\textbf{enonce.}
Comparateur - hysteresis : $V_{ref} = 2{,}42\,\mathrm{V}$, $V_{sat} = 15\,\mathrm{V}$, $R_1 = 1\,\mathrm{k\Omega}$, $R_2 = 150\,\mathrm{k\Omega}$.
(a) Exprimer $V_+$ en fonction de $V_{ref}$, $R_1$, $R_2$, $V_s$.
(b) Calculer $V_H$ et $V_B$.

\tcbline

\textbf{Corrige.}

\textbf{(a)} Diviseur de tension entre $V_s$ et $V_{ref}$ :
\[
  V_+ = \frac{R_2\,V_{ref} + R_1\,V_s}{R_1 + R_2} = \frac{150 \times 10^3 \times V_{ref} + 1 \times 10^3 \times V_s}{151 \times 10^3}
\]

\textbf{(b)} Seuil haut ($V_s = +V_{sat} = +15\,\mathrm{V}$) :
\[
  V_H = \frac{150\,000 \times 2{,}42 + 1\,000 \times 15}{151\,000} = \frac{363\,000 + 15\,000}{151\,000} = \frac{378\,000}{151\,000} \approx \mathbf{2{,}50\,\mathrm{V}}
\]

Seuil bas ($V_s = -V_{sat} = -15\,\mathrm{V}$) :
\[
  V_B = \frac{150\,000 \times 2{,}42 - 1\,000 \times 15}{151\,000} = \frac{363\,000 - 15\,000}{151\,000} = \frac{348\,000}{151\,000} \approx \mathbf{2{,}30\,\mathrm{V}}
\]

Hysteresis : $\Delta V = V_H - V_B = 2{,}50 - 2{,}30 = 0{,}20\,\mathrm{V}$ - protection contre les oscillations au voisinage du seuil.
\end{tcolorbox}

\vspace{4pt}

% -- QT3 -------------------------------------------------------------

\begin{tcolorbox}[questiontype,
  title={QT3 - Conversion tensionecourant - independance vis--vis de la charge \hfill\normalfont\footnotesize E2 disc. 2024 (E9)}
]
\textbf{enonce.}
Montage AOP en boucle fermee avec resistance de mesure $R_5 = 250\,\Omega$ dans la boucle. Tension en entree $V_4 \in [1\,;\,5]\,\mathrm{V}$. Montrer que le courant $I_s$ est independant de la charge $R_{ana}$ et calculer la plage.

\tcbline

\textbf{Corrige.}

Par la loi des neuds au neud A (entree $V_-$ de l'AOP, hypothese $I_- = 0$) :
\[
  \frac{V_A - V_4}{R_5} + \frac{V_A - V_s}{R_{ana}} = I_- = 0
\]
Avec $V_+ = V_- = V_A$ (contre-reaction), et en appliquant la regle d'entree $V_- = V_4$ :
\[
  I_s = \frac{V_4 - 0}{R_5} = \frac{V_4}{R_5}
\]
$I_s$ est bien \textbf{independant de $R_{ana}$} (la boucle de contre-reaction ajuste $V_s$ pour maintenir $\varepsilon = 0$ quelle que soit la charge).

\textbf{Plage de courant :}
\[
  V_4 = 1\,\mathrm{V} \;\Rightarrow\; I_s = \frac{1}{250} = 4\,\mathrm{mA}
  \qquad;\qquad
  V_4 = 5\,\mathrm{V} \;\Rightarrow\; I_s = \frac{5}{250} = 20\,\mathrm{mA}
\]
Standard industriel normalise 4-20\,mA.
\end{tcolorbox}

\end{tcolorbox}

% -----------------------------------------------------------------------
%  4. ATTENTES DU JURY
% -----------------------------------------------------------------------

\begin{tcolorbox}[
  jury,
  title={4. Attentes du jury - Erreurs recurrentes et points de vigilance},
]

\begin{enumerate}[label=\textbf{\arabic*.}, leftmargin=*, itemsep=5pt]

  \item \textbf{Verifier la contre-reaction avant d'appliquer $V_+ = V_-$.}
  En mode comparateur (boucle ouverte ou contre-reaction positive), $V_+ \neq V_-$ - ne pas appliquer l'egalite. Identifier d'abord le mode de fonctionnement.

  \item \textbf{Signe du gain inverseur.}
  $G = -R_2/R_1$ : le signe moins est fondamental. L'oubli conduit - conclure - une amplification non-inverseuse.

  \item \textbf{Courants d'entree AOP nuls.}
  $I_+ = I_- = 0$ : appliquer les lois de neuds uniquement sur les resistances externes. Ne jamais ecrire un courant passant par l'entree de l'AOP.

  \item \textbf{Theoreme de Millman meconnu.}
  Applique en 2024 (E4 : montage de decalage), il simplifie considerablement le calcul des neuds - plusieurs branches.

  \item \textbf{Hysteresis : ne pas inverser les seuils.}
  $V_H > V_B$ toujours. $V_H$ correspond - $V_s = +V_{sat}$, $V_B$ - $V_s = -V_{sat}$. Confondre les deux donne une hysteresis negative physiquement impossible.

  \item \textbf{Suiveur = adaptation d'impedance, pas amplification.}
  Gain = 1 mais impedance d'entree quasi infinie et impedance de sortie quasi nulle. Rele : isoler deux etages sans attenuation.

\end{enumerate}

\medskip
\begin{tcolorbox}[
  colback=colJuryFrame!12!colJury,
  colframe=colJuryFrame,
  left=4pt, right=4pt, top=2pt, bottom=2pt,
  fontupper=\small\itshape
]
Rapport jury 2024 : "La partie E (AOP) a ete abordee par 70\,\% des candidats et traitee correctement par seulement 47\,\%. Le jury observe un manque de maetrise des montages complexes (conversion T/I, decalage). Les calculs numeriques doivent etre encadres avec des unites."\\
Rapport jury 2024 : "Connaetre les trois conditions de l'AOP ideal et les appliquer systematiquement sans les redemontrer ($V_+ = V_-$, $I_+ = I_- = 0$). Justifier le montage choisi parmi les montages standards."
\end{tcolorbox}

\end{tcolorbox}

% -----------------------------------------------------------------------
%  PIED - PReSENCE DANS LES SUJETS
% -----------------------------------------------------------------------

\begin{tcolorbox}[
  enhanced, breakable,
  colback=colGrayLight,
  colframe=colSeparator,
  left=4pt, right=4pt, top=3pt, bottom=3pt,
  before skip=4pt, after skip=0pt,
  fonttitle=\small\bfseries,
  title={Presence dans les sujets analyses (2013-2026)},
]
\begin{center}
\footnotesize\renewcommand{\arraystretch}{1.15}
\begin{tabular}{lccccccccccccc}
\toprule
\textbf{Annee} & 13 & 14 & 15 & 16 & 17 & 18 & 19 & 20 & 21 & 22 & 23 & 24 & 25 \\
\midrule
E1 (commune)   & - & - & - & - & - & - & - & - & - & - & - & $\checkmark$ & - \\
E2 discipl. IE & - & - & - & - & - & - & $\checkmark$ & - & - & $\checkmark$ & $\checkmark$ & $\checkmark$ & - \\
E2 appl. comm. & - & - & - & - & - & - & - & - & - & - & - & - & - \\
\bottomrule
\end{tabular}
\end{center}
\smallskip
\footnotesize\textbf{Sous-themes frequents :}
montage inverseur/non-inverseur (6/13) - montage soustracteur (4/13) - comparateur - hysteresis (3/13) - conversion tension-courant 4-20\,mA (2/13) - integrateur/correcteur PI (3/13)

\medskip
\textbf{Ordres de grandeur (reperes rapides).}
\begin{itemize}[leftmargin=1.3em]
  \item Reseau BT FR : tension phase-neutre typique \SI{230}{\volt}, tension phase-phase typique \SI{400}{\volt}, frequence \SI{50}{\hertz}.
  \item Courants d'etude frequents en sujet : de quelques \si{\ampere} a quelques dizaines de \si{\ampere}.
  \item En conversion d'energie, un rendement attendu reste generalement entre 0{,}8 et 0{,}98 selon la technologie.
\end{itemize}

\medskip
\textbf{Synthese express (a refaire sans support).}
\begin{enumerate}[leftmargin=1.4em]
  \item Poser les hypotheses du modele et les conventions de signe.
  \item Ecrire 2 a 4 relations cles avec unites.
  \item Realiser un calcul numerique complet et verifier l'ordre de grandeur.
  \item Justifier le resultat par un argument technique de niveau CAPET.
\end{enumerate}

\end{tcolorbox}

\CAPETENDFILE









